\documentclass[11pt,a4paper]{article}
\usepackage[utf8]{inputenc}
\usepackage{amsmath,amssymb}
\usepackage{geometry}
\geometry{margin=2.5cm}
\usepackage{listings}
\usepackage{xcolor}
\usepackage{hyperref}

\lstset{
  language=C++,
  basicstyle=\ttfamily\small,
  keywordstyle=\color{blue}\bfseries,
  commentstyle=\color{green!60!black},
  stringstyle=\color{red!70!black},
  numbers=left,
  numberstyle=\tiny\color{gray},
  numbersep=8pt,
  frame=single,
  breaklines=true,
  tabsize=4,
  showstringspaces=false
}

\title{IOI 1997 -- Mars Explorer}
\author{Solution}
\date{}

\begin{document}
\maketitle

\section{Problem Statement}

A robot navigates an $n \times n$ grid where some cells are obstacles.
Starting at $(1,1)$, it must reach $(n,n)$, return to $(1,1)$, and go to $(n,n)$ again ---
three traversals total. On forward trips the robot moves only right or down;
on the return trip it moves only left or up.

Maximize the number of \textbf{distinct} non-obstacle cells visited across all
three traversals.

\medskip
\noindent\textbf{Constraints:} $n \le 30$.

\section{Solution Approach}

\subsection{Reduction to Three Simultaneous Forward Paths}

A return trip from $(n,n)$ to $(1,1)$ using left/up moves is equivalent to a forward
trip from $(1,1)$ to $(n,n)$ using right/down moves (simply reverse the path).
Thus the problem reduces to finding three forward paths from $(1,1)$ to $(n,n)$
that collectively visit the maximum number of distinct cells.

\subsection{DP on Three Simultaneous Paths}

We advance all three paths in lockstep. At step $t$ (where $0 \le t \le 2(n-1)$),
each path occupies a cell $(r_i, c_i)$ with $r_i + c_i = t + 2$. Since $c_i$ is
determined by $t$ and $r_i$, the state is $(t, r_1, r_2, r_3)$.

\medskip
\noindent\textbf{Value:} Maximum distinct cells visited by positions $1, \ldots, t$.

\medskip
\noindent\textbf{Transition:} Each path moves right ($r_i$ unchanged) or down ($r_i + 1$),
giving $2^3 = 8$ combinations per state.

\subsection{Symmetry Optimization}

Without loss of generality, enforce $r_1 \le r_2 \le r_3$. This reduces the number of
states by a factor of up to 6.

\subsection{Counting Distinct Cells}

At each step, the three paths occupy cells $(r_1, c_1)$, $(r_2, c_2)$, $(r_3, c_3)$.
Since they share the same $t$, distinct cells correspond to distinct $r$-values:
positions with the same $r_i$ also share $c_i$. The number of new distinct cells
is $|\{r_1, r_2, r_3\}|$ (provided all cells are non-obstacles).

\section{C++ Solution}

\begin{lstlisting}
#include <cstdio>
#include <cstring>
#include <algorithm>
using namespace std;

int n;
int grid[32][32]; // 1 = free, 0 = obstacle
int dp[62][32][32][32]; // dp[t][r1][r2][r3]

// Returns number of distinct free cells, or -1 if any cell is invalid
int countDistinct(int t, int r1, int r2, int r3) {
    int c1 = t - r1 + 2, c2 = t - r2 + 2, c3 = t - r3 + 2;
    if (r1 < 1 || r1 > n || c1 < 1 || c1 > n || !grid[r1][c1]) return -1;
    if (r2 < 1 || r2 > n || c2 < 1 || c2 > n || !grid[r2][c2]) return -1;
    if (r3 < 1 || r3 > n || c3 < 1 || c3 > n || !grid[r3][c3]) return -1;

    int cnt = 1;
    if (r2 != r1) cnt++;
    if (r3 != r1 && r3 != r2) cnt++;
    return cnt;
}

int main() {
    scanf("%d", &n);
    for (int i = 1; i <= n; i++)
        for (int j = 1; j <= n; j++)
            scanf("%d", &grid[i][j]);

    memset(dp, -1, sizeof(dp));

    int maxT = 2 * (n - 1);

    // Base case: t = 0, all three paths at (1,1)
    if (grid[1][1])
        dp[0][1][1][1] = 1;

    for (int t = 0; t < maxT; t++) {
        for (int r1 = 1; r1 <= n; r1++) {
            int c1 = t - r1 + 2;
            if (c1 < 1 || c1 > n) continue;
            for (int r2 = r1; r2 <= n; r2++) {
                int c2 = t - r2 + 2;
                if (c2 < 1 || c2 > n) continue;
                for (int r3 = r2; r3 <= n; r3++) {
                    int c3 = t - r3 + 2;
                    if (c3 < 1 || c3 > n) continue;
                    if (dp[t][r1][r2][r3] < 0) continue;

                    int cur = dp[t][r1][r2][r3];

                    // Try all 8 move combinations (right or down for each path)
                    for (int d1 = 0; d1 <= 1; d1++)
                    for (int d2 = 0; d2 <= 1; d2++)
                    for (int d3 = 0; d3 <= 1; d3++) {
                        int nr1 = r1 + d1, nr2 = r2 + d2, nr3 = r3 + d3;

                        // Sort to maintain r1 <= r2 <= r3 invariant
                        int sr[3] = {nr1, nr2, nr3};
                        sort(sr, sr + 3);

                        int add = countDistinct(t + 1, sr[0], sr[1], sr[2]);
                        if (add < 0) continue;

                        int& ref = dp[t+1][sr[0]][sr[1]][sr[2]];
                        ref = max(ref, cur + add);
                    }
                }
            }
        }
    }

    printf("%d\n", dp[maxT][n][n][n]);
    return 0;
}
\end{lstlisting}

\section{Complexity Analysis}

\begin{itemize}
  \item \textbf{Time complexity:} $O(n^4)$.
        There are $2(n-1) + 1$ time steps and $O(n^3/6)$ ordered triples per step
        (after the symmetry reduction). Each state tries 8 transitions. Total:
        $O(n \cdot n^3 / 6 \cdot 8) = O(n^4)$.
        For $n = 30$: roughly $30 \cdot 30^3 / 6 \cdot 8 \approx 1{,}080{,}000$ operations.
  \item \textbf{Space complexity:} $O(n^4)$ for the DP table. For $n = 30$:
        $60 \times 30^3 \approx 1{,}620{,}000$ entries, each a 4-byte integer.
\end{itemize}

\end{document}
